\section{Degree-Preserving Edge Intervention}
\label{sec:edge_intervention}

\subsection{Paired Intervention Design}

We separately test whether GCN uses edge organization beyond degree on Cora, CiteSeer, and PubMed. This is auxiliary evidence about graph signal on the three citation networks; the largest diagnostic-loss settings involve different graphs and selected architectures. For each dataset and seed, verified double-edge swaps randomize the graph while preserving the complete node set, edge count, and degree of every node. The procedure completes five successful swaps per original edge, with randomization seed equal to the split seed plus 100,000. The runner stores both edge lists and recomputes the simple-graph, self-loop, duplicate-edge, per-node-degree, connectivity, and homophily checks before training.

The original and randomized conditions use the same node features, labels, split, GCN architecture, and four-trial hyperparameter grid. They select hyperparameters independently by validation accuracy and evaluate the selected test state once. The paired estimand is randomized minus original test accuracy for the same dataset, split, and seed. This design isolates sensitivity to edge organization beyond the degree sequence, but it does not isolate a single changed property: swaps jointly alter class mixing, connectivity patterns, motifs, and higher-order neighborhoods.

\subsection{Results}

\begin{table}[t]
\centering
\caption{Degree-Preserving Edge-Randomization Results. Differences and confidence intervals are percentage points. Each interval is a percentile-bootstrap interval for the mean of ten paired seed differences within that dataset, using 10,000 resamples. It describes variation under the fixed dataset and specified randomization procedure; the three-dataset interval is reported separately in the text.}
\label{tab:degree_intervention}
\small
\begin{tabular}{@{}lcccc@{}}
\toprule
Dataset & Original (\%) & Randomized (\%) & Difference & 95\% CI \\
\midrule
Cora & 88.0 & 42.4 & $-45.5$ & $[-47.2,-44.1]$ \\
CiteSeer & 77.2 & 47.4 & $-29.9$ & $[-30.6,-29.1]$ \\
PubMed & 87.3 & 69.9 & $-17.3$ & $[-17.8,-16.9]$ \\
\bottomrule
\end{tabular}
\end{table}

Table~\ref{tab:degree_intervention} summarizes the intervention; all 30 paired seed differences are negative. The seed-level two-sided sign-flip tests yield Holm-adjusted $p=0.005859$ for each dataset. The equal-weight mean across the three dataset effects is $-30.9$ points, with a descriptive dataset-bootstrap interval of $[-45.5,-17.3]$. However, three dataset clusters provide only eight possible sign patterns; the exact dataset-cluster sign-flip test gives $p=0.25$.

GCN accuracy is sensitive to edge organization beyond the degree sequence on these three citation networks. Because randomization jointly changes several structural properties, the effect cannot be attributed to homophily alone. This auxiliary evidence separates the use of graph structure by a trained model from the utility of a diagnostic for choosing that model.
